% Gerado por scripts/gerar-dicionario-dados.py.
\resizebox{0.95\linewidth}{!}{%
\begin{tikzpicture}[entity/.style={draw,rounded corners,fill=white,align=center,text width=3.4cm,minimum height=1.25cm,inner sep=4pt,font=\small},fk/.style={-{Latex},semithick,draw=black!70}]
\node[entity] (n6) at (0,-8) {\texttt{class\_\allowbreak{}exercise\_\allowbreak{}lists}\\7 atributos};
\node[entity] (n3) at (0,-4) {\texttt{class\_\allowbreak{}members}\\3 atributos};
\node[entity] (n4) at (5,-4) {\texttt{classes}\\8 atributos};
\node[entity] (n10) at (10,-12) {\texttt{exercise\_\allowbreak{}list\_\allowbreak{}items}\\4 atributos};
\node[entity] (n9) at (0,-12) {\texttt{exercise\_\allowbreak{}lists}\\8 atributos};
\node[entity] (n5) at (10,-4) {\texttt{exercises}\\9 atributos};
\node[entity] (n2) at (10,0) {\texttt{languages}\\13 atributos};
\node[entity] (n0) at (0,0) {\texttt{organizations}\\3 atributos};
\node[entity] (n7) at (5,-8) {\texttt{submissions}\\11 atributos};
\node[entity] (n8) at (10,-8) {\texttt{test\_\allowbreak{}cases}\\6 atributos};
\node[entity] (n1) at (5,0) {\texttt{users}\\9 atributos};
\begin{scope}[on background layer]
\draw[fk] (n6) -- (n9);
\draw[fk] (n6) -- (n4);
\draw[fk] (n3) -- (n4);
\draw[fk] (n3) -- (n1);
\draw[fk] (n4.north west) -- (n0.south east);
\draw[fk] (n4) -- (n1);
\draw[fk] (n10) -- (n9);
\draw[fk] (n10.east) -- (12.1,-12) -- (12.1,-4) -- (n5.east);
\draw[fk] (n9.east) -- (2.5,-12) -- (2.5,-1.8) -- (n1.south west);
\draw[fk] (n9.south) -- (0,-13.5) -- (12.5,-13.5) -- (12.5,0) -- (n2.east);
\draw[fk] (n5) -- (n1);
\draw[fk] (n5) -- (n2);
\draw[fk] (n2) -- (n1);
\draw[fk] (n2) to[loop above,min distance=18mm] (n2);
\draw[fk] (n7) -- (n5);
\draw[fk] (n7) -- (n9);
\draw[fk] (n7) -- (n4);
\draw[fk] (n7.north east) -- (7.5,-6.5) -- (7.5,-1.8) -- (n1.south east);
\draw[fk] (n8) -- (n5);
\draw[fk] (n1) -- (n0);
\draw[fk] (n1) to[bend left=25] (n2);
\end{scope}
\node[font=\small,align=left,text width=14cm,anchor=north west,inner sep=0pt] at (-1.8,-14.1) {Cada seta parte da tabela que contém a FK e aponta para a tabela referenciada.\\Setas em sentidos opostos entre usuários e linguagens representam propriedade e linguagem ativa. A autorreferência de linguagens registra sua origem por clonagem.};
\end{tikzpicture}}
